% Technical appendix. Mechanics and diagnostics moved out of the main text so
% Section 3 states what the instrument establishes rather than how it is built.
% Unlimited appendix pages are permitted; the main-text limit is nine.
% No literal numbers: every figure is a macro from numbers.tex.

\appendix

\section{Extractor mechanics}
\label{app:extractor}

\paragraph{Locating a declaration block.} The release ships no source
positions, so a declaration's block is found textually: the declaration
keyword under any attribute or modifier prefix, extending to the next line that
begins a top-level command. The terminator set is deliberately broad, because
the two failure modes are not symmetric. Over-matching truncates a block and
loses tokens at its tail, while under-matching swallows the following
declaration and invents them. Failures to locate are dominated by auto-named
instances, which have no textual declaration site.

\paragraph{Tokenization.} Comments and string literals are blanked to spaces,
preserving token boundaries; Lean's block comments nest, so they are scanned
rather than matched by regular expression. Identifier tokens are then taken in
Lean's own shape, dot-composed, with a leading dot excluded. The \texttt{.mp}
in \texttt{h.mp} is projection on a local, not a reference to a root-level
\texttt{mp}. A short keyword list drops tactic names before
resolution. It is a precision aid rather than the safety net. Every surviving
token must still resolve against the declaration set, and the list is
deliberately short, since an over-long one silently excludes real premises that share a name
with a tactic. Section~\ref{sec:limitations}, L3b, reports the one inconsistency
we left in it rather than extending it after seeing which names flattered the
result.

\paragraph{Resolution.} The rule is three steps, first hit wins: the token is a
fully-qualified declaration name; or it completes to exactly one declaration
under a namespace prefix in scope; or it resolves to nothing and is kept aside
as a bare tail. Bare tails never enter the premise set, for the reason
Section~\ref{sec:extraction} gives.

\paragraph{Namespace prefixes.} The prefixes a bare token may complete under
are the proper prefixes of the declaration's own name, together with the
\texttt{open} clauses appearing in its file. Lean scopes \texttt{open} to
sections and this does not, so the rule is over-permissive, and it can
complete a token under a namespace not actually open at that point. The cost
is precision and never recall, and it is carried in the reference line of
Section~\ref{sec:reference-line}.

\section{Residual below the declaration reference}
\label{app:residual}

The \residualTargets{} mathematical targets that the declaration reference
point finds and the extractor misses decompose into three classes.
\residualResolutionPct{}\% is a definitional gap: the reference point
substring-matches premise names while the extractor resolves identifiers, and a
resolved identifier is a stricter object than a matched substring, so closing
that gap entirely would mean abandoning resolution and buying retention at the
cost of precision. \residualKeywordPct{}\% is the over-broad keyword list of
Appendix~\ref{app:extractor}, which suppresses names that are simultaneously
tactic names and library declarations; we report its cost rather than tuning
the list after seeing which names it flatters. The remaining
\residualCommentPct{}\% occurs only inside comments or string literals, which
the extractor strips by design. The classes are identified rather than merely
counted, so each is a statement about a mechanism.

Per-class extraction efficiency (achieved retention against the declaration
reference point, within each tactic class) ranges from
\efficiencyByClassMin{} to \efficiencyByClassMax{}, with the corpus-level
figure \efficiencyReachable{} inside that range. Section~\ref{sec:limitations},
L1, states what this does and does not license about the machine corpus.

\section{Instrument reference line: diagnostics}
\label{app:reference}

The precision term is carried as an open item in
Section~\ref{sec:reference-line}, and the diagnostics behind that decision are
these. Length-controlled on the tactic-step axis the term falls to
\spuriousStepAxis{}, against \spuriousCharAxis{} on the character axis, a
near-collapse that read alone would be good news about the instrument. By a
criterion registered before the run, a control whose collapsed cells carry a
large share of the machine mix is uninformative rather than clean: here they
carry \spuriousCollapsedMachineMixPct{}\% of it, dominated by the arithmetic
stratum of Section~\ref{sec:concentration}. The character axis fails the
companion criterion of near-disjoint supports, its overlap coefficient being
\charOverlapSample{}. One diagnostic points somewhere real. Across the two
axes the imputed human rate is stationary while only the machine rate moves,
which is the signature of a compositional correction rather than of two
extrapolations converging, and that makes the step figure the more credible of
the two without making it something to freeze on.

\section{The tactic-step proxy}
\label{app:steps}

Counting tactic steps needs an instrument and the obvious one under-records.
The published tactic lists capture only main-path tactics; everything after
\texttt{<;>}, inside \texttt{induction \dots\ with} case branches, and in
nested \texttt{(by simp \dots)} term positions is missing. That is measured
rather than assumed. Against a string matcher over the same declarations the
disagreement is almost entirely one-directional, the matcher finding tactics
the lists omit and rarely the reverse. Both corpora are therefore counted with
a single fixed tactic vocabulary and one matcher. It is a size proxy, not a
taxonomy, and a tactic missing from the vocabulary costs sensitivity
symmetrically. Section~\ref{sec:limitations}, L2, states the consequence for
every stratum and mix in the paper.

\section{Structural breadth by design}
\label{app:breadth}

\begin{table}[h]
\centering
\caption{Breadth per tactic step, machine minus human, at three scopes, with
the share of the machine corpus each design retains. The sign does not move;
the magnitude moves by a factor of four, which is why
Section~\ref{sec:breadth} quotes a direction and not a magnitude.}
\label{tab:breadth}
\small
\begin{tabular}{lcc}
\toprule
Design & Gap & Machine corpus retained \\
\midrule
Full corpora   & \breadthGapFull{}    & all \\
Common support & \breadthGapSupport{} & \supportRetentionMachine{} \\
Matched 1:1    & \breadthGapMatched{} & \supportRetentionMachine{} \\
\bottomrule
\end{tabular}
\end{table}

Two things keep breadth secondary. The community partition is heavily skewed,
with most communities tiny and a few holding much of the graph, so the unit is
coarse before any comparison is made; and a quantity whose magnitude depends this strongly on the
design cannot carry weight on its own. We report it nonetheless because the
metric was defined before we looked, and omitting an unstable result would be
selection rather than caution. Breadth also measures how far apart a
population's premises lie within the library, rather than how much of the
library that population reaches, so it is not a second reading of coverage.

The excluded region under the last two designs is the set of cells with no
human counterpart, overwhelmingly the arithmetic stratum of
Section~\ref{sec:concentration}. An earlier version of this comparison did
move in sign. Restricting to a length-overlap region gave a positive gap
against the full corpora's negative one. That region was cut on character
length, which
Section~\ref{sec:length-axis} shows to be a formatting control rather than a
length control, and it retained \charOverlapMachineRetainedPct{}\% of the
machine corpus. On the wrap-invariant axis the sign is stable at all three
scopes, so the flip is an artifact of the rejected axis and not a finding.

\section{Limitations L2, L3, L3d and L3e}
\label{app:limitations}

\paragraph{L2. Instrument asymmetries, and their direction.} Human tactic use
comes from the publishers' extracted lists and machine tactic use has no
equivalent, so both sides are cut with the string matcher of
Section~\ref{sec:length-axis} and Appendix~\ref{app:steps}; the two instruments
disagree materially and in one direction, and one false positive is kept rather
than removed after the fact, because the tactic list was frozen before the
machine corpus was touched. On the extraction side, two artifacts that
resolved on the human side only (a declaration's own name and the
\texttt{lemma} keyword) are removed before any reported figure, at a large
cost to our own headline, and three further asymmetries are left in place: \texttt{symm} and
\texttt{trans} resolve \symmTransHuman{} times on the human side against
\symmTransMachine{}, prefix completion supplies far more premises per proof to
the human side because every machine proof carries the same fixed namespace
header, and coverage is counted over both hub layers where the mathematical
layer alone would give \vocabRatioMathOnly{}$\times$ rather than
\vocabRatio{}$\times$. Every uncorrected asymmetry we have found narrows the
reported gap, so Section~\ref{sec:coverage}'s figures are floors. The
corrections we applied cost us, and the ones we declined would have paid us.
Separately, no aggregate centrality figure appears without the hub layers
separated and we never infer shallowness from centrality;
Section~\ref{sec:dag-layer} reports where we broke that rule ourselves.

\paragraph{L3. Extraction cannot recover the explicit set in full.} Sections
\ref{sec:is-explicit} and \ref{sec:retention} characterize this rather than
bound it away. The whole-file reference point is \ceilingWholeFile{},
which no per-declaration instrument beats, and it is a loose-matcher figure
rather than an attributable bound. The reported vocabularies are
instrument-filtered.

\paragraph{L3b/L3c. The asymmetries that remain all shrink the gap.} Two
artifacts that resolved on the human side only, a declaration's own name
and the \texttt{lemma} keyword, are removed before any reported figure, at a
large cost to our own headline. Three further asymmetries are
left in place and all run the same way. \texttt{symm} and \texttt{trans} are
tactic names that are also mathlib theorems, resolving \symmTransHuman{} times
on the human side against
\symmTransMachine{}; we leave
them rather than extend a keyword list after seeing which names flatter the
result. Prefix completion supplies far more premises per proof to the human
side,
since every machine proof carries the same fixed namespace header while a
mathlib declaration sits in its own. That is legitimate resolution rather than
contamination. And coverage is counted over both hub layers, where the
mathematical layer alone would give \vocabRatioMathOnly{}$\times$ rather than
\vocabRatio{}$\times$. Every uncorrected asymmetry we have found narrows the
reported gap, so Section~\ref{sec:coverage}'s figures are floors. The
corrections we applied cost us, and the ones we declined would have paid us.

\paragraph{L3d. The reference points are matcher-based.} The quantities of
Table~\ref{tab:retention} are cut with a loose name matcher: a match is
evidence that a name occurs within the scope, not that the premise was used
there, and we do not measure how often it credits an incidental occurrence.

\paragraph{L3e. Our reading of the flag is not validated against elaborated
terms.} The semantics established in Section~\ref{sec:is-explicit} come from
the release's own parser and the documented behaviour of the tool upstream of
it, not from elaborating proofs ourselves, which we do not do anywhere in this
work.

\section{A quarter of the machine vocabulary is not mathematics}
\label{sec:layer-composition}

The retention figures of Sections~\ref{sec:retention} and
\ref{sec:reference-line} are quoted on the mathematical layer, using a
partition of the declaration set by provenance rather than by namespace
prefix: Lean core constants, instances, coercions, kernel-level kinds and
compiler-generated names are infrastructure, and the rest is mathematical
content. Provenance is what makes the partition usable. A prefix rule misses
every root-level elaborator constant, and it over-captures genuine mathlib
results whose names happen to sit under an infrastructure namespace.
Applied to the premise vocabularies themselves rather than to retention
targets, the same partition says something about the corpora.

The human vocabulary is \humanVocabInfraPct{}\% infrastructure:
\humanVocabInfra{} declarations against \humanVocabMath{} mathematical ones.
The machine vocabulary is \machineVocabInfraPct{}\% infrastructure,
\machineVocabInfra{} against \machineVocabMath{}, a factor of
\infraWeightRatio{} more infrastructure-weighted. Both vocabularies place
completely, with no declaration left unclassified on either side.

The machine vocabulary therefore differs from the human one in composition,
not only in size. A quarter of it is the plumbing of
Section~\ref{sec:is-explicit} (proof-irrelevance and rewriting machinery that
every Lean proof drags in and no author selects), so the mathematical working
vocabulary of the machine
corpus is \machineVocabMath{} declarations. That is the figure to hold in mind
next to mathlib's \humanVocabMath{}, since it is the set of mathematical
results the released corpus actually reaches for, as opposed to the set it
cites.

This is a membership count and nothing more. Each premise is classified once,
by provenance, and the two sides are compared by how many fall in each class.
There is no aggregation over a centrality score, so Li et al.'s Finding~3 does
not bind here and we draw no inference about mathematical depth from it. A
vocabulary weighted toward infrastructure is not thereby a shallower one.

This composition sits next to the concentration reported in
Section~\ref{sec:concentration}, where most of the machine corpus is short
automation calls. A proof that discharges its goal in one such call still
elaborates into a term, and that term carries the rewriting and
proof-irrelevance machinery with it, so the two measurements are consistent
with one picture of the corpus. We report them as two measurements rather than
as one explained by the other. For evaluation the consequence is direct. A
vocabulary that is \machineVocabInfraPct{}\% infrastructure and one that is
\humanVocabInfraPct{}\% are different objects, and only the second is evidence
of mathematical range, a distinction no correctness-based metric records.

The split also settles a basis question left open in
Section~\ref{sec:coverage}. Coverage is reported over both layers. Counting it
on the mathematical layer alone would give \vocabRatioMathOnly{}$\times$
rather than \vocabRatio{}$\times$, with a correspondingly higher human-only
fraction, because the machine vocabulary is more infrastructure-weighted and
so loses proportionally more under the restriction. We report the both-layer
basis, which is the weaker of the two for our own hypothesis, and print the
other here. Where two bases are defensible, the one to keep is the one that
raises the bar the result has to clear.

\section{A characteristic concentration}
\label{sec:concentration}

Coverage counts each premise once, which makes it a set rather than a rate and
is why Section~\ref{sec:coverage} reports it with no per-proof standardization;
it also makes coverage blind to usage. Counting the same premises once per
occurrence gives a different picture. The machine corpus resolves
\machineOccurrences{} premise occurrences from its \machineVocab{} distinct
premises, against \humanOccurrences{} from \humanVocab{} on the human side, and
it concentrates on a stratum essentially absent from mathlib: proofs closing in
at most two steps with an arithmetic-flavoured tactic set are
\degenerateMachineSharePct{}\% of the machine corpus and
\degenerateHumanSharePct{}\% of mathlib, \degenerateHumanCount{} of
\humanProofs{} declarations. We report that as a result rather than as sample
attrition. It is the largest cell of the machine corpus, it has no human
counterpart, and it is why the common-support design sets roughly a quarter of
the machine corpus aside. Usage concentrates to match. \texttt{sq\_nonneg}
alone accounts for \sqNonnegOccurrencePct{}\% of machine premise occurrences
and the twenty-five most frequent for \topTwentyFiveOccurrencePct{}\%, both
descriptive and post-hoc, and both not regenerable from the released artifact
(Section~\ref{sec:limitations}, L6). A machine proof resolves
\machinePremisesPerProof{} premises on average against
\humanPremisesPerProof{} for a human declaration. Coverage and concentration
are independent and neither substitutes for the other. Heavy reuse of a few
lemmas is ordinary in competition mathematics, and coverage alone would leave
open that a shared vocabulary is used comparably.

\section{\texttt{dag\_layer}: the reading under which it is correct}
\label{app:daglayer}

Section~\ref{sec:dag-layer} records the column as a centrality proxy and the
error in reading it as depth as ours. The attribute may be behaving as its
producers intended: if edges run from a declaration to its dependencies and the
computation seeds at in-degree-zero nodes, a column that puts uncited leaves at
layer zero and \texttt{Eq.refl} at the top is exactly what the definition
yields. Nothing here is a defect report against the release.

A second hazard is worth naming for other consumers. \layerTruncatedPct{}\% of
nodes carry a partial value that never finished propagating and is
indistinguishable in the published data from a completed one, so a consumer who
filters or bins on this column is mixing two populations without being told.

\section{What the coverage quantities do and do not read}
\label{app:construct}

These quantities read no edges. The objection is fair and we state it first.
Computing $J$ consults only membership, so coverage is a measure on node sets. The
graph supplies the definition of a node (which names are declarations, and
which layer each belongs to) rather than the arithmetic. The section headings retain the phrase \emph{structural
coverage}, and the reader should take it as naming the pair of measures defined
here rather than as a claim about the coverage arithmetic. Coverage is the
node-set family; breadth is the edge-derived one.

Two reporting conventions follow from that. We print the intersection count
beside every Jaccard so each index is recomputable from its own row, and we
give the human-only quantity as both a count and a fraction, because the two
vocabularies differ in size by orders of magnitude and a ratio alone would hide
which side is large. Coverage and breadth also answer different questions and
can move independently: coverage asks how much of the library a population
reaches, breadth how spread out that reach is, and a corpus may touch few
premises scattered across many communities or many premises within one.

The choice is deliberate rather than a convenience. Any edge-reading statistic
over these subgraphs requires aggregation: a centrality, a path length, a depth.
Aggregation over this graph is what Li et al.'s Finding~3 shows to be
uninterpretable unless the two hub layers are separated first, and
Section~\ref{sec:dag-layer} reports a second instance of the same trap on a
different quantity. A membership count needs no aggregation, so it needs no such
split to be read, and it survives controls that reverse or destroy the
edge-reading quantities.
